先在带边界的紧商空间 $\widetilde\Omega=\Omega\times(\mathbb R/\mathbb Z)^m$ 上工作，将数据沿新增的周期方向作常值延拓。在四槽基 $(1,\mathbf i,\mathbf j,\mathbf k)$ 中，采用指标 $0,1,2,3$，原周期方向 $\Lambda$ 位于第 $0$ 槽，新周期方向 $\mathbb Z^m$ 位于第 $3$ 槽。重组实坐标，可将 $\widetilde\Omega$ 等同于环面 $\mathbb R^{2m}/(\Lambda\times\mathbb Z^m)$（第 $0,3$ 槽）与 $V$（第 $1,2$ 槽）的乘积。在降回原空间之前，$G,H$ 均表示完整的四槽梯度与 Hessian。记 $\mathcal H(D^2\psi)=H(\psi)$，并用 $\widehat\xi$ 表示实余向量 $\xi$ 按四槽编码所得的四元数列向量。引理~\ref{ms:lemma-3-1} 给出 $\mathcal H(\xi\xi^{\mathsf t})=\widehat\xi\widehat\xi^*$、$\operatorname{tr}H(\psi)=\Delta\psi$ 及在四元数酉坐标变换下的协变性。

\emph{比较和延拓。} 从已知解 $u=r$ 出发，按以下路径连接到所求的密度
\[
 g_t=(1-t)\det H(B_0+r)+tg,\qquad 0\le t\le1.
\]
在解 $\psi=B_0+u$ 处，令 $Q=H(\psi)$ 与 $S=\operatorname{tr}Q^{-1}$。由引理~\ref{ms:lemma-3-2}，对数行列式的导数及其实主符号为
\[
 Lw=Q^{-1}:H(w),\qquad \sigma_L(\xi)=\widehat\xi^*Q^{-1}\widehat\xi .
\]
因此 $L$ 在任一固定光滑解处一致椭圆，条件是在紧闭包上满足 $Q>0$。若两个势函数满足 $Q_i=H(\psi_i)>0$，对相应矩阵作插值得
\[
 \log\det Q_1-\log\det Q_2=
 \left(\int_0^1((1-\theta)Q_2+\theta Q_1)^{-1}\,\mathrm d\theta\right):H(\psi_1-\psi_2).
\]
当 Hessian 直至边界均正定时，平均后的实算子一致椭圆。由最大值原理得到四槽情形下的比较结论，并推出唯一性。同一论证适用于有界子域或带边界的紧商空间。

对每个 $k\ge0$ 和 $0<\alpha<1$，线性化 Dirichlet 算子 $L:C_0^{k+2,\alpha}\to C^{k,\alpha}$ 为同构，其中下标表示零边界迹。\cite{r-31}, Theorems 6.14 and 6.19 中的局部线性 Schauder 估计和 Dirichlet 理论适用于平直商空间的有限坐标覆盖。使用连续性方法，从 Laplace 算子变形到 $L$。在此过程中，定义函数 $r$ 始终是严格的障碍函数，因为 $H(r)\ge cI$。最大值原理给出此线性插值的一致 $C^0$ 估计，Schauder 估计给出相应的高阶界，且起始的 Dirichlet Laplace 算子可逆。现在对非线性方程应用隐函数定理，得到可解参数 $t$ 的集合为开集。

在此乘积空间上使用四槽 Hessian，引理~\ref{ms:dirichlet-c2} 为路径上的每个光滑解给出两个障碍函数以及一致的完整 $C^2$ 估计。

\emph{正则性、降回原空间与参数依赖。} $C^2$ 界和 $\det Q=g_t\ge c_0>0$ 给出 $cI\le Q\le CI$。因此实 Hessian 均位于紧凸集
\[
 \mathcal K=\{M:|M|\le C_1,\ cI\le\mathcal H(M)\le CI\}.
\]
上。在 $\mathcal H(M)>0$ 上，光滑函数 $F(M)=\log\det\mathcal H(M)$ 由引理~\ref{ms:lemma-3-2} 可知为凹函数。对 $M\in\mathcal K$，有
\[
 DF(M)[\xi\xi^{\mathsf t}]=
 \widehat\xi^*\mathcal H(M)^{-1}\widehat\xi
 \in[C^{-1}|\xi|^2,c^{-1}|\xi|^2].
\]
这些正是 $\mathcal K$ 上的实一致椭圆性与凹性假设，适用于 \cite{r-31}, Theorems 17.14 and $17.26'$ 的内部及边界估计，即 hypotheses (i)$'$ 和 (ii)$'$。为使算子定义在整个实矩阵空间上，将 $F$ 乘以一个光滑截断函数，令其支撑于满足 $\mathcal H(M)>0$ 的位置，且在 $\mathcal K$ 的一个邻域内等于一。估计所用的方程为 $F(D^2\psi)-\log g_t=0$，光滑边界值为 $B_0$。在有限坐标覆盖中应用这些局部估计，可得一致的 $C^{2,\alpha}$ 界，其中 $\alpha>0$ 取某个值。在可解参数的极限处，取一个在 $C^{2,\beta}$ 中收敛的子列，其中 $0<\beta<\alpha$。正定性与方程在极限下保持。通过 Schauder 估计逐步提高正则性，得到直至边界的光滑性：拉平边界后，利用给定的 Dirichlet 数据作切向求导，再由一致椭圆性从方程恢复纯法向二阶导数。这就证明了闭性，所以路径到达 $g$。

沿新增周期坐标的平移保持数据及区域不变。由 $\widetilde\Omega$ 上解的唯一性，解在所有这些平移下不变，因而可降到 $\Omega$。物理空间中的任意解都可提升，由此得到唯一性。通过同一提升也可得到物理子域上的比较原理。对于光滑族，将提升后的问题拉回固定域。其系数及边值光滑依赖于参数，且线性化算子与 Dirichlet 同构共轭。在每个 $C_0^{k+2,\alpha}$ 中应用隐函数定理，即得直到边界的光滑参数依赖。用于选取 $b,b'$ 的严格行列式不等式在该族中局部保持成立。
